\begin{theorem}[Group-derived equivalence for all closed group insertions]
    \label{thm:peps_graph_inserted_semiregular_factory}
    \lean{TNLean.PEPS.exists_isometric_minimalSemiRegular_graphInsertedStates}
    \leanok
    \uses{thm:peps_graph_minimal_semiregular_equivalence,
        thm:peps_graph_inserted_semiregular_equivalence}
    Let $G$ be finite and let $\Gamma$ be a finite simple graph. There are
    matrix representations $D_i$ of dimensions $d_i$, a unitary
    semi-regular representation $U=\bigoplus_i D_i$, and one full-domain
    physical bond isometry whose product $\mathcal T$ satisfies
    \begin{align}
        \mathcal T\Psi_{U,\Theta,u} & =\Psi_{L,I,u}
                                    &                             & \text{for every }u:E(\Gamma)\to G, \\
        \langle\mathcal T\psi,\mathcal T\varphi\rangle
                                    & =\langle\psi,\varphi\rangle
                                    &                             & \text{for every }\psi,\varphi.
    \end{align}
    Here $\Theta=\bigoplus_i d_i^{1/4}I_{d_i}$, and the states are the
    actual closed averaging-site contractions with the representation
    insertion $U(u_e)$, respectively $L(u_e)$, on each oriented bond.
    Every displayed character has multiplicity one in $U$. Moreover,
    \begin{align}
        \sum_i d_i\leq\dim\sigma
    \end{align}
    for every finite-dimensional complex semi-regular representation
    $\sigma$, without requiring $\sigma$ to be unitary.
    The representation and the physical isometry are chosen before all
    insertion assignments and all physical vectors. No Fourier,
    coefficient or Gram identity is supplied as a hypothesis.
    This is an auxiliary inserted-state consequence of
    \cite[Section~7]{Schuch2010PEPS} on closed finite simple graphs;
    it does not assert equivalence of arbitrary open-boundary images
    or an isometry for arbitrary matrices mixing sectors.
\end{theorem}

\begin{proof}\leanok
    Reuse the group-derived representation witnesses from the closed-state
    existence theorem, including their unitary, semi-regular and universal
    minimality properties. Apply the uniform inserted-state theorem to
    those same witnesses. Its full-domain bond isometry is chosen before
    all edge insertions, and its product preserves every inner product.
\end{proof}
